\documentclass[11pt]{article}

%========================================================
% PAGE GEOMETRY AND BASIC SETTINGS
%========================================================
\usepackage[margin=1in]{geometry}
\usepackage{microtype}
\usepackage{setspace}
\setstretch{1.15}

%========================================================
% MATH PACKAGES AND SYMBOLS
%========================================================
\usepackage{amsmath, amssymb, amsthm, amsfonts, mathtools, bm}
\usepackage{dsfont}
\usepackage{mathrsfs}
\usepackage{bbm}
\usepackage{stmaryrd}
\usepackage{esint}
\usepackage{colonequals}

%========================================================
% TABLES, FIGURES, AND CAPTIONS
%========================================================
\usepackage{booktabs}
\usepackage{threeparttable}
\usepackage{float}
\usepackage{graphicx}
\usepackage[font=small,labelfont=bf]{caption}

%========================================================
% ENUMERATION AND TEXT UTILITIES
%========================================================
\usepackage{enumitem}
\setlist[itemize]{leftmargin=1.25em, itemsep=3pt, topsep=3pt}
\setlist[enumerate]{leftmargin=1.25em, itemsep=3pt, topsep=3pt}

\usepackage{xcolor}
\usepackage{url}
\usepackage[colorlinks=true,linkcolor=black,citecolor=black,urlcolor=blue]{hyperref}

%========================================================
% NUMBERING AND THEOREM STYLES
%========================================================
\numberwithin{equation}{section}

\newtheorem{theorem}{Theorem}[section]
\newtheorem{lemma}[theorem]{Lemma}
\newtheorem{proposition}[theorem]{Proposition}
\newtheorem{corollary}[theorem]{Corollary}

\newtheorem{assumption}{Assumption}[section]
\newtheorem*{assumption*}{Assumption}

\newlist{assumpenum}{enumerate}{1}
\setlist[assumpenum,1]{label=\textup{(\alph*)}, ref=\theassumption\textup{(\alph*)}, leftmargin=1.5em}

\theoremstyle{definition}
\newtheorem{definition}[theorem]{Definition}

\theoremstyle{remark}
\newtheorem{remark}[theorem]{Remark}

%========================================================
% CUSTOM MATH OPERATORS AND SHORTCUTS
%========================================================

% Probability and expectation
\DeclareMathOperator{\E}{\mathbb{E}}
\DeclareMathOperator{\Var}{Var}
\DeclareMathOperator{\Cov}{Cov}
\DeclareMathOperator{\Prob}{\mathbb{P}}
\DeclareMathOperator{\plim}{plim}

% Optimization and algebra
\DeclareMathOperator*{\argmin}{argmin}
\DeclareMathOperator*{\argmax}{argmax}
\DeclareMathOperator{\diag}{diag}
\DeclareMathOperator{\rank}{rank}
\DeclareMathOperator{\spec}{spec}
\DeclareMathOperator{\Sym}{Sym}
\DeclareMathOperator{\Id}{I}

% Differential operators
\newcommand{\dd}{\mathrm{d}}
\newcommand{\pd}{\partial}
\newcommand{\grad}{\nabla}

% Useful short macros
\newcommand{\1}{\mathds{1}}
\newcommand{\R}{\mathbb{R}}
\newcommand{\N}{\mathbb{N}}

% Bracket sizes
\newcommand{\paren}[1]{\left(#1\right)}
\newcommand{\brac}[1]{\left[#1\right]}
\newcommand{\bracec}[1]{\left\{#1\right\}}

% Common bold symbols
\newcommand{\bA}{\mathbf{A}}
\newcommand{\bB}{\mathbf{B}}
\newcommand{\bC}{\mathbf{C}}
\newcommand{\bD}{\mathbf{D}}
\newcommand{\bE}{\mathbf{E}}
\newcommand{\bH}{\mathbf{H}}
\newcommand{\bI}{\mathbf{I}}
\newcommand{\bK}{\mathbf{K}}
\newcommand{\bL}{\mathbf{L}}
\newcommand{\bM}{\mathbf{M}}
\newcommand{\bP}{\mathbf{P}}
\newcommand{\bR}{\mathbf{R}}
\newcommand{\bw}{\mathbf{w}}
\newcommand{\bp}{\mathbf{p}}
\newcommand{\bone}{\mathbf{1}}

% Common bold Greek symbols
\newcommand{\balpha}{\boldsymbol{\alpha}}
\newcommand{\bbeta}{\boldsymbol{\beta}}
\newcommand{\btheta}{\boldsymbol{\theta}}
\newcommand{\btau}{\boldsymbol{\tau}}
\newcommand{\bpsi}{\boldsymbol{\psi}}
\newcommand{\bPi}{\boldsymbol{\Pi}}

%========================================================
% PROJECT-SPECIFIC MACROS
%========================================================
% Add note-specific operators and notation here.

%========================================================
% DOCUMENT APPEARANCE AND SECTION STYLING
%========================================================
\usepackage{titlesec}
\titleformat{\section}{\large\bfseries}{\thesection.}{0.5em}{}
\titleformat{\subsection}{\normalsize\bfseries}{\thesubsection.}{0.5em}{}
\titleformat{\paragraph}[runin]{\bfseries}{\theparagraph}{0.5em}{}
\titlespacing*{\section}{0pt}{1.2em plus 0.3em minus 0.1em}{0.8em}
\titlespacing*{\subsection}{0pt}{1em plus 0.2em minus 0.1em}{0.5em}
\titlespacing*{\paragraph}{0pt}{0.8em plus 0.2em minus 0.1em}{0.5em}

%========================================================
% HYPERREF SETUP
%========================================================
\hypersetup{
  colorlinks=true,
  linkcolor=black,
  citecolor=black,
  urlcolor=blue,
  pdfauthor={},
  pdftitle={Research Note},
  pdfsubject={Summary and technical appendix},
  pdfkeywords={economics, theory, empirical, computational}
}

%========================================================
% TITLE AND AUTHOR INFO
%========================================================
\title{\textbf{Research Note Title}\\[0.5em]
\large Subtitle Describing the Main Object or Result}
\author{}
\date{\today}

\begin{document}
\maketitle

\section{Summary of Main Results}\label{sec:summary}

\paragraph{Goal.}
State the economic question and the object this note derives or summarizes.

\paragraph{Main result.}
State the main result in words before giving notation-heavy details.

\paragraph{Status.}
State whether the result is proved, numerically verified, empirically estimated, conjectural, or a maintained approximation.

\section{Setup and Notation}\label{sec:setup}

Define the economic environment, primitives, observables, policy variables, and equilibrium objects.

\section{Main Results}\label{sec:main-results}

\begin{proposition}[Main result]\label{prop:main}
Under the maintained assumptions, the economic object of interest admits the representation
\begin{equation}\label{eq:main-representation}
  Y = A + B.
\end{equation}
\end{proposition}

\begin{proof}
The full derivation is in Appendix~\ref{app:derivations}. The key step is to decompose the total effect into the direct component \(A\) and the equilibrium adjustment \(B\).
\end{proof}

\section{Interpretation}\label{sec:interpretation}

Explain the economic content of the result. Identify the direct effect, equilibrium response, incidence, substitution, selection, or identification channel.

\section{Checks, Limitations, and Open Issues}\label{sec:checks}

\paragraph{Checks.}
List benchmark cases, sign checks, units, limiting cases, and computational or empirical diagnostics.

\paragraph{Limitations.}
State what the note does not yet prove, estimate, or verify.

\appendix

\section{Derivations}\label{app:derivations}

\subsection{Deriving the main representation}\label{app:derive-main}

\paragraph{Economic object.}
We derive \eqref{eq:main-representation}, which decomposes the object \(Y\) into a direct component and an adjustment component.

\paragraph{Primitive relationship.}
Start from
\begin{equation}\label{eq:primitive}
  Y = f(x,z),
\end{equation}
where \(x\) is the direct margin and \(z\) is an endogenous adjustment.

\paragraph{Line-by-line derivation.}
Differentiating \eqref{eq:primitive} with respect to a policy shifter \(\tau\) gives
\begin{align}
  \frac{\dd Y}{\dd \tau}
  &= \frac{\pd f(x,z)}{\pd x}\frac{\dd x}{\dd \tau}
   + \frac{\pd f(x,z)}{\pd z}\frac{\dd z}{\dd \tau} \notag\\
  &= A + B. \label{eq:main-derivative}
\end{align}
The first equality applies the chain rule. The second defines \(A\) as the direct effect through \(x\) and \(B\) as the adjustment effect through \(z\).

\paragraph{Economic interpretation.}
Equation \eqref{eq:main-derivative} separates the mechanical effect of the policy from the endogenous equilibrium response.

\end{document}
